\documentclass[../main]{subfiles}
\begin{document}

\chapter{Primer Principio de la Termodinámica}

\section{Primer Principio de la Termodinámica}

    \subsection{Introducción al Primer Principio de la Termodinámica}

El Primer Principio de la Termodinámica es una de las leyes fundamentales de la física que describe cómo se conserva la energía en los sistemas termodinámicos. Establece que la energía no puede ser creada ni destruida, solo transformada de una forma a otra. Esta ley es crucial para comprender el funcionamiento de los motores térmicos, refrigeradores y muchas otras máquinas que utilizan energía térmica\footnote{La energía termica es la energía almacenada en un cuerpo}.

    \info{Primer Principio de la Termodinámica}{
El Primer Principio de la Termodinámica se expresa matemáticamente como:

\begin{equation}
\Delta U = Q - W
\end{equation}

Donde:

\begin{itemize}
    \item $\Delta U$: Cambio en la energía interna del sistema [J].
    \item $Q$: Calor añadido al sistema [J].
    \item $W$: Trabajo realizado por el sistema [J].
\end{itemize}

El calor $Q$ es la energía transferida debido a una diferencia de temperatura, mientras que el trabajo $W$ es la energía transferida cuando una fuerza actúa a través de una distancia. En términos simples, si se añade calor a un sistema o se realiza trabajo sobre él, su energía interna aumenta. Si el sistema realiza trabajo o cede calor, su energía interna disminuye.
    }

    \info{Aplicaciones del Primer Principio de la Termodinámica}{
El Primer Principio de la Termodinámica se aplica en múltiples situaciones prácticas, como en el análisis de ciclos termodinámicos. Un ejemplo clásico es el ciclo de Carnot, que es un modelo idealizado de un motor térmico. Este ciclo describe cómo el calor puede transformarse en trabajo mecánico de manera eficiente.

Otra aplicación importante es en los procesos de calefacción y refrigeración. En un refrigerador, por ejemplo, el objetivo es extraer calor del interior y liberarlo al exterior. El trabajo necesario para mover este calor está relacionado con el cambio en la energía interna del sistema y el calor transferido.
    }

    \info{Ecuación de Estado de los Gases Ideales}{
Para muchos problemas prácticos, especialmente a nivel de secundaria, es útil combinar el Primer Principio de la Termodinámica con la ecuación de estado de los gases ideales, que es:

\begin{equation}
PV = nRT
\end{equation}

Donde:

\begin{itemize}
    \item $P$: Presión del gas [Pa].
    \item $V$: Volumen del gas [m³].
    \item $n$: Número de moles del gas [mol].
    \item $R$: Constante de los gases ideales [$8.314 \, \text{J/(mol·K)}$].
    \item $T$: Temperatura absoluta del gas [K].
\end{itemize}

Esta ecuación relaciona la presión, el volumen y la temperatura de un gas ideal, y es fundamental para resolver problemas de termodinámica en los que se involucran estos parámetros.
    }

\actividad{Responder el siguiente cuestionario}
\begin{enumerate}
    \item ¿Qué establece el Primer Principio de la Termodinámica?
    \item ¿Cómo se define la energía interna de un sistema?
    \item ¿Qué es el calor y cómo se transfiere?
    \item ¿Qué es el trabajo en términos de termodinámica?
    \item ¿Cómo se expresa matemáticamente el Primer Principio de la Termodinámica?
    \item ¿Qué aplicaciones prácticas tiene el Primer Principio de la Termodinámica?
    \item ¿Qué es un ciclo de Carnot?
    \item ¿Cómo se relacionan el calor, el trabajo y la energía interna en un ciclo termodinámico?
    \item ¿Qué es la ecuación de estado de los gases ideales y para qué se utiliza?
    \item ¿Cómo se relacionan la presión, el volumen y la temperatura en un gas ideal?
\end{enumerate}

\actividad{Resolver los siguientes problemas}
\begin{enumerate}
    \item A un sistema se le añaden $500 \, \text{J}$ de calor y realiza un trabajo de $200 \, \text{J}$. ¿Cuál es el cambio en su energía interna?
    \item Un gas ideal a $300 \, \text{K}$ y $1 \, \text{atm}$ ocupa un volumen de $0.1 \, \text{m}^3$. Si se duplica la presión a temperatura constante, ¿cuál será su nuevo volumen?
    \item Un sistema realiza un trabajo de $150 \, \text{J}$ y libera $250 \, \text{J}$ de calor al entorno. ¿Cuál es el cambio en su energía interna?
    \item En un proceso isocórico, se añaden $400 \, \text{J}$ de calor a un gas ideal. ¿Cuál es el cambio en su energía interna?
    \item Un gas ideal realiza un trabajo de $100 \, \text{J}$ cuando su volumen cambia de $0.5 \, \text{m}^3$ a $1 \, \text{m}^3$ a presión constante. ¿Cuál es su presión?
    \item Un sistema recibe $600 \, \text{J}$ de calor y su energía interna aumenta en $400 \, \text{J}$. ¿Cuánto trabajo realiza el sistema?
    \item Un gas ideal a $273 \, \text{K}$ ocupa un volumen de $22.4 \, \text{L}$. Si se duplica la temperatura a presión constante, ¿cuál será su nuevo volumen?
    \item Un sistema tiene un cambio en su energía interna de $-300 \, \text{J}$ y realiza un trabajo de $100 \, \text{J}$. ¿Cuánto calor ha intercambiado con su entorno?
    \item Un gas ideal a $500 \, \text{K}$ y $2 \, \text{atm}$ ocupa un volumen de $0.05 \, \text{m}^3$. Si se reduce la temperatura a la mitad a presión constante, ¿cuál será su nuevo volumen?
    \item Un sistema absorbe $800 \, \text{J}$ de calor y realiza un trabajo de $300 \, \text{J}$. Si su cambio en la energía interna es de $200 \, \text{J}$, ¿cuál es el signo del calor transferido?
\end{enumerate}

\end{document}


\end{document}
